% =====================================================================
% Section 1 — Введение  (REVTeX 4-2, twocolumn, run-in headings)
% Owner:  ALL (initial draft by N. Novitsky)
%
% Style notes (mirroring пример/main.tex + project rules):
%   - section headings use the project macro \sect{Title}{label},
%     which renders as a BOLD inline run-in heading
%   - no enumerate / itemize lists in the body
%   - minimal use of \textbf and \textit (only for headings + first-time
%     foreign terms)
%   - centered display equations, numbered on the right via \eqlab{}
% =====================================================================

\sect{Введение}{introduction}
%
Автоматический аккомпанемент относится к числу старейших и наиболее
требовательных задач анализа музыкальной информации. Солист ожидает,
что оркестр будет следовать за его темпом, динамикой и артикуляцией в
реальном времени, однако подавляющее большинство учащихся
консерваторий не имеют возможности регулярно репетировать концертные
произведения с живым ансамблем. Существующие записи-минусовки
навязывают исполнителю фиксированный темп и игнорируют
индивидуальное rubato, формирующее художественную интерпретацию. В
рамках мастерской MusicTech Центрального Университета (Т-Банк)
разрабатывается приложение виртуального оркестра, которое
воспринимает игру солиста через MIDI или микрофон и синхронно
воспроизводит оркестровую партию.

Сложность задачи определяется тремя факторами. Во-первых,
исполнитель допускает значительные временные отклонения от
номинального темпа партитуры, причём знак и величина отклонения не
известны заранее и могут изменяться от такта к такту. Во-вторых,
живое исполнение содержит ошибки, пропуски, повторы и орнаменты, не
предусмотренные нотной записью. В-третьих, аккомпанемент должен быть
выдан с задержкой не более $20$\,мс, что соответствует порогу
восприятия рассинхронизации в музыкальном ансамбле.

\sect{Постановка задачи}{problem}
%
Формально задача отслеживания партитуры состоит в построении
отображения
%
\begin{equation}
\varphi : \{1, \dots, M\} \longrightarrow \{1, \dots, N\},
\qquad j \mapsto \varphi(j),
\eqlab{alignment}
\end{equation}
%
сопоставляющего каждому событию исполнения с номером~$j$ позицию
$\varphi(j)$ в партитуре, состоящей из~$N$ нот. Здесь и далее $M$~---
общее число наблюдаемых событий исполнения, $N$~--- число нот
партитуры. Поскольку исполнитель движется по партитуре вперёд (с
возможными пропусками, но без возвращения назад), отображение
$\varphi$ должно удовлетворять условию монотонности
%
\begin{equation}
j_1 < j_2 \quad\Longrightarrow\quad \varphi(j_1) \le \varphi(j_2).
\eqlab{monotonicity}
\end{equation}
%
Каузальность задачи означает, что значение $\varphi(j)$ необходимо
вычислить, имея в распоряжении только наблюдения $q_1,\dots,q_j$ и
не зная будущих событий $q_{j+1},\dots,q_M$.

\sect{Существующие подходы}{approaches}
%
Подходы к решению этой задачи можно условно разделить на три группы.
Первую группу образуют алгоритмы динамического программирования
(Dynamic Time Warping и его онлайн-вариант On-line Time Warping):
они дают глобально или локально оптимальное выравнивание, но не
располагают вероятностной моделью неопределённости. Ко второй группе
относятся скрытые марковские модели, естественно описывающие
неопределённость через апостериорное распределение по позициям, но
требующие ручной настройки переходных вероятностей. Третью группу
составляют нейросетевые методы, обеспечивающие высокую точность на
сложных полифонических данных, однако требующие значительных
вычислительных ресурсов и больших обучающих выборок.

\sect{Вклад настоящей работы}{contributions}
%
В настоящей работе предлагается комбинированная схема, объединяющая
сильные стороны указанных групп. Во-первых, представлено единое
математическое описание моделей HMM и OLTW для монофонического входа
в формате MIDI с числовыми примерами, воспроизводимыми на
синтетическом датасете. Во-вторых, описан конвейер обработки в
реальном времени, включающий очистку входных событий, объединение
двух трекеров (HMM как основной, OLTW как резервный) и потоковый
вывод оркестрового аккомпанемента с явным распределением задержек по
этапам. В-третьих, предложен компактный свёрточный модуль для оценки
апостериорного распределения по высотам нот при работе с
аудиовходом, заменяющий гауссову функцию эмиссии HMM, а также
систематическая классификация граничных случаев и правил их
обработки. В-четвёртых, в качестве основного методологического
вклада статьи предлагается модуль на основе обучения с подкреплением,
предсказывающий изменения темпа на горизонте $200$--$500$\,мс и
передающий результат базовому трекеру; этот модуль снижает
наблюдаемую задержку аккомпанемента без потери точности выравнивания.

Структура изложения соответствует приведённому перечню вкладов и
последовательно раскрывается в следующих разделах работы.
